\documentclass[11pt,a4paper]{article}
\usepackage[a4paper,left=2.3cm,right=2.3cm,top=2cm,bottom=2.2cm]{geometry}
\usepackage{fontspec}
\setmainfont{PT_Sans-Web}[Path=fonts/, Extension=.ttf, UprightFont=ptsans_*-Regular, BoldFont=ptsans_*-Bold, ItalicFont=ptsans_*-Italic, BoldItalicFont=ptsans_*-BoldItalic]
\setsansfont{PT_Sans-Web}[Path=fonts/, Extension=.ttf, UprightFont=ptsans_*-Regular, BoldFont=ptsans_*-Bold, ItalicFont=ptsans_*-Italic, BoldItalicFont=ptsans_*-BoldItalic]
\newfontfamily\slab{ZillaSlab}[Path=fonts/, Extension=.ttf, UprightFont=zillaslab_*-SemiBold, BoldFont=zillaslab_*-SemiBold]
\usepackage{xcolor}
\definecolor{accent}{HTML}{003366}
\definecolor{magenta}{HTML}{CC0066}
\definecolor{muted}{HTML}{64748B}
\color[HTML]{243649}
\usepackage[hidelinks]{hyperref}
\hypersetup{pdftitle={Gonzalo Llamosas García, currículum}, pdfauthor={Gonzalo Llamosas García}, pdflang={es}}
\usepackage{titlesec}
\usepackage{fancyhdr}
\usepackage{calc}
\setlength{\parindent}{0pt}
\setlength{\parskip}{0pt}
\titleformat{\section}{\slab\color{accent}\Large}{{\color{magenta}\rule[-0.15em]{3pt}{1em}}\hspace{0.5em}}{0pt}{}
\titlespacing*{\section}{0pt}{14pt}{6pt}
\newcommand{\entry}[2]{\begin{minipage}[t]{2.6cm}\sffamily\small\textcolor{muted}{#1}\end{minipage}%
  \begin{minipage}[t]{\linewidth-2.6cm}\raggedright #2\end{minipage}\par\vspace{4pt}}
\pagestyle{fancy}\fancyhf{}\renewcommand{\headrulewidth}{0pt}
\fancyfoot[L]{\sffamily\footnotesize\color{muted}Gonzalo Llamosas García, currículum actualizado en octubre de 2026}
\fancyfoot[R]{\sffamily\footnotesize\color{muted}\thepage}
\begin{document}
{\slab\color{accent}\fontsize{26}{30}\selectfont Gonzalo Llamosas García}\par\vspace{6pt}
{\sffamily\large Profesor Ayudante Doctor de Hacienda Pública}\par\vspace{2pt}
{\sffamily Departamento de Economía Aplicada (Hacienda Pública, Política Económica y Economía Política), Universidad de Málaga}\par\vspace{6pt}
{\sffamily\small\color{muted}Despacho 2402, Facultad de Ciencias Económicas y Empresariales, Campus de El Ejido, 29071 Málaga\\
\href{mailto:gonzalo.llamosas@uma.es}{gonzalo.llamosas@uma.es} \quad \href{https://gonzalollamosas.com}{gonzalollamosas.com} \quad ORCID \href{https://orcid.org/0000-0002-4407-2845}{0000-0002-4407-2845}}\par
\vspace{4pt}{\sffamily\small Acreditado por ANECA como Profesor Contratado Doctor}\par
\section{Trayectoria académica}
\entry{2025–}{Profesor Ayudante Doctor de Hacienda Pública\\{\small\color{muted}Universidad de Málaga, Departamento de Economía Aplicada}}
\entry{2024–2025}{Profesor Ayudante Doctor\\{\small\color{muted}Universidad de Cantabria}}
\entry{2022–2024}{Profesor sustituto interino\\{\small\color{muted}Universidad de Málaga}}
\entry{2023–2024}{Profesor\\{\small\color{muted}Universidad Europea del Atlántico}}
\entry{2020–2022}{Profesor\\{\small\color{muted}Universidad Loyola Andalucía}}
\entry{2015–2019}{Investigador predoctoral en Economía Aplicada\\{\small\color{muted}Universidad de Cantabria}}

\section{Formación}
\entry{2015–2019}{Doctor en Economía, Universidad de Cantabria (programa de doctorado interuniversitario «Economía: Instrumentos del Análisis Económico», con las universidades de Oviedo y del País Vasco)\\{\small\color{muted}Tesis: Essays on Citizens’ Economic Behavior in Services of Public Interest: A European Union Perspective.}}
\entry{2012–2013}{Máster en Investigación en Economía, Universidad de Valladolid}
\entry{2008–2012}{Licenciatura en Economía, Universidad de Cantabria}

\section{Artículos en evaluación}
\entry{}{Decomposing the Debt Response to Fiscal Consolidation: Evidence from EU-27 Local Projections, 1996–2024\\{\small\color{muted}con Cristina Mazas Pérez-Oleaga y José M. Domínguez-Martínez}}

\section{Documentos de trabajo}
\entry{}{Event-study estimators when adoption responds to past outcomes}
\entry{}{Does Biological Ageing Foreshadow Financial Fragility? Evidence from Dried-Blood-Spot Biomarkers and Household Wealth in Europe}

\section{Publicaciones}
\entry{2025}{Llamosas-García, G. y Mazas Pérez-Oleaga, C. «What Works in Financial Education? Experimental Evidence on Program Impact». \emph{Journal of Behavioral and Experimental Economics}, vol. 119, 102401. \href{https://doi.org/10.1016/j.socec.2025.102401}{\color{accent}DOI} (acceso abierto)}
\entry{2023}{Lezcano Lastra, A., Llamosas-García, G. y López-Cagigas, A. «Predicting the Tide of the Pandemic: An In-Depth Analysis of Forecasting Models for COVID-19 in Cantabria». \emph{BEIO. Boletín de Estadística e Investigación Operativa}, vol. 39, n.º 2, pp. 36–49. \href{https://portaldelainvestigacion.uma.es/documentos/6501a724df0d06753d5bfb47}{\color{accent}Enlace}}
\entry{2020}{Clifton, J., Díaz-Fuentes, D. y Llamosas-García, G. «A Behavioral Perspective on Saving Decisions. Empirical Evidence for Policymakers in the European Union». \emph{Global Policy}, vol. 11, supl. 1, pp. 62–72. \href{https://doi.org/10.1111/1758-5899.12751}{\color{accent}DOI} (acceso abierto)}
\entry{2020}{Clifton, J., Díaz-Fuentes, D. y Llamosas-García, G. «ICT-Enabled Co-Production of Public Services: Barriers and Enablers. A Systematic Review». \emph{Information Polity}, vol. 25, n.º 1, pp. 25–48. \href{https://doi.org/10.3233/IP-190122}{\color{accent}DOI} (lectura gratuita)}

\section{Trabajos en curso}
\entry{}{Sustainability for Whom? The Distributional Costs of Statutory Retirement-Age Reforms in Older Europe}
\entry{}{Corporate Tax Structures in the Age of AI: Evidence from World Bank Data}
\entry{}{Popular Discontent and Recall Elections}

\section{Proyectos de investigación}
\entry{2020–2021}{\emph{EU Finance and Institutions: New Social and Environmental Challenges}, Módulo Jean Monnet, Comisión Europea. Investigador.}
\entry{2017–2020}{\emph{Europe, Finance and Institutions}, Módulo Jean Monnet, Comisión Europea. Investigador.}
\entry{2016–2019}{\emph{CITADEL: Empowering Citizens to Transform European Public Administrations}, Horizonte 2020, Unión Europea. Investigador predoctoral.}
\entry{2015–2017}{\emph{Global Banks Programme}, Santander Financial Institute (SANFI). Ayudante de investigación.}

\section{Estancias de investigación}
\entry{2021–2022}{Loyola Behavioral Lab, Universidad Loyola Andalucía (de octubre de 2021 a abril de 2022), con Pablo Brañas-Garza}
\entry{2019–2020}{Estancia posdoctoral como profesor visitante, Universidad Internacional Iberoamericana (UNIB), Puerto Rico (de octubre de 2019 a octubre de 2020)}
\entry{2017}{Bangor Business School, Universidad de Bangor, Reino Unido (abril y mayo de 2017), con Phil Molyneux}
\entry{2016}{Departamento de Administración de Empresas, Universidad de Valladolid (abril y mayo de 2016), con Eleuterio Vallelado}

\section{Premios}
\entry{2016}{Premio a la mejor comunicación, I Workshop en Gobierno, Incentivos y Riesgo en la Banca (Sección de Economía Financiera de ACEDE y Santander Financial Institute), Santander, por «Loss Aversion in Saving Behavior in Europe»}
\entry{2012}{II Premio a la Comunicación por «Desarrollo Económico en Argentina», XIV Premio del Ilustre Colegio de Economistas de Cantabria}

\section{Labor editorial y evaluación de artículos}
\entry{2016–2018}{Editor asistente, Journal of Economic Policy Reform}
\entry{Revisor}{Journal of Behavioral and Experimental Economics; Public Finance Review; Revista Galega de Economía}

\section{Comunicaciones en congresos}
\entry{2023}{\emph{Evaluando la trayectoria experimental de la educación financiera. Una revisión sistemática}. VI Congreso de Educación Financiera de Edufinet: «La educación financiera en un contexto de transición hacia la sostenibilidad», Málaga}
\entry{2021}{\emph{Una evaluación de modelos de pronóstico para casos COVID-19 en Cantabria}. XXI Jornadas de Estadística de las Comunidades Autónomas (JECAS 2021), Canarias}
\entry{2019}{\emph{Behavioral Perspective on Saving Behavior}. Banca March Workshop on Contemporary Issues in Banking, Madrid}
\entry{2019}{\emph{Explaining citizens’ coproduction through ICTs in health: a multi-level model analysis}. IIAS Study Group on Co-production of Public Services, KU Leuven}
\entry{2017–2018}{Asambleas generales del proyecto CITADEL: La Haya, Riga, Bari, Bruselas y Amberes}
\entry{2017}{\emph{Co-creation in financial services: a semantic analysis}. Wolpertinger Conference de la European Association of University Teachers of Banking and Finance, Santander}
\entry{2016}{\emph{Loss Aversion in Saving Behavior: Evidence from the ECB’s European Household Finance and Consumption Survey}. XVIII Reunión de Economía Mundial, Universidad de Alcalá}
\entry{2016}{\emph{Loss Aversion in Saving Behavior in Europe: Evidence from the Household Finance and Consumption Survey}. I Workshop en Gobierno, Incentivos y Riesgo en la Banca, Santander}
\entry{2015}{\emph{¿Tomas riesgos en tus inversiones? Tus características socioeconómicas pueden hacerte perder dinero}. II Jornadas de Jóvenes Investigadores de Cantabria, Santander}

\section{Docencia}
\entry{2025–}{Universidad de Málaga\\{\small\color{muted}Hacienda Pública; Hacienda Pública Española; Sistemas Fiscales; dirección de Trabajos Fin de Grado}}
\entry{2022–2024}{Universidad de Málaga\\{\small\color{muted}Sistemas Fiscales; Microeconomía; Economía Contemporánea de Asia Oriental}}
\entry{2024–2025}{Universidad de Cantabria\\{\small\color{muted}Producción y Organización Industrial; Economía y Administración de Empresas; Empresas (grados de ingeniería); Dirección y Gestión de Empresas (Máster Universitario en Ingeniería Industrial)}}
\entry{2023–2024}{Universidad Europea del Atlántico\\{\small\color{muted}Business Organisation y Market Research (docencia en inglés); dirección de Trabajos Fin de Grado}}
\entry{2022–2023}{Universidad de León\\{\small\color{muted}Mercados, Medios e Instituciones Financieras (Grado en Administración y Dirección de Empresas)}}
\entry{2020–2022}{Universidad Loyola Andalucía\\{\small\color{muted}Economía Internacional; Política Económica; Aplicaciones II: Estudio de casos}}

\section{Otra experiencia profesional}
\entry{2021–2023}{Técnico de Estadísticas Económicas, Gobierno de Cantabria\\{\small\color{muted}Análisis estadístico de datos administrativos, modelos de predicción e informes para las políticas de sanidad, educación y finanzas públicas.}}

\section{Idiomas e informática}
\entry{Idiomas}{Español (nativo), inglés (C1), francés (básico)}
\entry{Informática}{R, Python, Stata, SQL, LaTeX}

\end{document}
